\documentclass[12pt]{article}
\usepackage{amsmath, amssymb}
\usepackage{enumitem}

\title{LECTURE SCRIBE – CSE400 (Lecture 18)}
\date{}

\begin{document}

\maketitle

\section*{Topics}
\begin{itemize}
    \item Marginal PMF
    \item Correlation
    \item Covariance
    \item Pearson’s Coefficient
    \item Joint Gaussian Random Variables
\end{itemize}

\section{Marginal PMF}

\subsection*{Definition}

\textbf{Key Idea (from slide):} \\
The marginal PMF of one variable is obtained by \textbf{summing the joint PMF over all possible values of the other variable}.

\subsection*{Marginal PMF from Joint PMF Table}

\textbf{Step-by-step reasoning:}

Let joint PMF be $P(X = x, Y = y)$

\subsubsection*{To find marginal of $X$:}

\begin{enumerate}
    \item Fix a value of $x$
    \item Sum over all possible values of $y$
\end{enumerate}

\[
P_X(x) = \sum_y P(X=x, Y=y)
\]

\subsubsection*{To find marginal of $Y$:}

\begin{enumerate}
    \item Fix a value of $y$
    \item Sum over all possible values of $x$
\end{enumerate}

\[
P_Y(y) = \sum_x P(X=x, Y=y)
\]

\subsection*{Example (as given in slides)}

\textbf{Given:} \\
Joint PMF

\textbf{Solution steps (as shown):}

\begin{enumerate}
    \item Identify all values of joint PMF
    \item For each $x$, sum across all $y$
    \item For each $y$, sum across all $x$
\end{enumerate}

\textbf{Final conclusion:} \\
Marginal PMF is obtained by \textbf{systematically summing joint probabilities along rows or columns}.

\section{Correlation (Raw Correlation)}

\subsection*{Definition}

A basic way to measure joint behavior is through the \textbf{product moment}:

\[
R_{XY} = E[XY]
\]

\subsection*{Interpretation}

\begin{itemize}
    \item Represents the \textbf{average value of the product $(XY)$}
    \item Called a \textbf{raw correlation measure}
\end{itemize}

\subsection*{Important Observations (from slides)}

\begin{itemize}
    \item Depends on \textbf{scale of $X$ and $Y$}
    \item Computed \textbf{about the origin (not about means)}
    \item Mixes:
    \begin{itemize}
        \item Mean values
        \item Spread
        \item Dependence
    \end{itemize}
\end{itemize}

\subsection*{Special Case}

\[
R_{XY} = 0 \Rightarrow \text{orthogonal about the origin}
\]

\subsection*{Transition Idea}

To remove the effect of means:
\begin{itemize}
    \item \textbf{Center the variables}
    \item Leads to \textbf{Covariance}
\end{itemize}

\section{Covariance}

\subsection*{Definition}

\[
\text{Cov}(X,Y) = E[(X - E[X])(Y - E[Y])]
\]

\subsection*{Step-by-Step Interpretation}

\begin{enumerate}
    \item Subtract mean from each variable:
    \begin{itemize}
        \item $(X - E[X])$
        \item $(Y - E[Y])$
    \end{itemize}
    \item Multiply deviations:
    \[
    (X - E[X])(Y - E[Y])
    \]
    \item Take expectation:
    \[
    E[(X - E[X])(Y - E[Y])]
    \]
\end{enumerate}

\subsection*{Meaning (from slides)}

\begin{itemize}
    \item Measures how \textbf{centered quantities vary together}
    \item Captures whether deviations occur:
    \begin{itemize}
        \item In same direction
        \item In opposite direction
    \end{itemize}
\end{itemize}

\section{Algebraic Form of Covariance}

\subsection*{Start with definition:}

\[
\text{Cov}(X,Y) = E[(X - E[X])(Y - E[Y])]
\]

\subsection*{Expand step-by-step:}

\[
= E[XY - X E[Y] - Y E[X] + E[X]E[Y]]
\]

\subsection*{Apply expectation:}

\[
= E[XY] - E[X]E[Y] - E[Y]E[X] + E[X]E[Y]
\]

\subsection*{Simplify:}

\[
\boxed{\text{Cov}(X,Y) = E[XY] - E[X]E[Y]}
\]

\section{Properties of Covariance}

\subsection*{1. Symmetry}
\[
\text{Cov}(X,Y) = \text{Cov}(Y,X)
\]

\subsection*{2. Variance as Special Case}
\[
\text{Var}(X) = \text{Cov}(X,X)
\]

\subsection*{3. Linearity with Respect to Constants}

Covariance is linear when constants are involved.

\subsection*{4. Independence Property}

If $X$ and $Y$ are independent:
\[
\text{Cov}(X,Y) = 0
\]

\textbf{Important conclusion (from slide):}
\begin{itemize}
    \item Zero covariance $\Rightarrow$ \textbf{uncorrelated}
    \item Does \textbf{not imply independence}
\end{itemize}

\section{Covariance Example 1}

\textbf{Given:} \\
Expressions for variables

\textbf{Solution Steps (exact slide flow):}

\begin{enumerate}
    \item Use covariance formula:
    \[
    \text{Cov}(X,Y) = E[(X - E[X])(Y - E[Y])]
    \]
    \item Substitute given expressions into formula
    \item Expand expression
    \item Compute expectations
    \item Substitute values into covariance formula
    \item Simplify
\end{enumerate}

\textbf{Final Result:} \\
Obtained after substituting and simplifying (as shown in slides).

\section{Pearson’s Correlation Coefficient}

\subsection*{Definition}

\begin{itemize}
    \item Measures \textbf{effect of change in one variable when the other changes}
\end{itemize}

\subsection*{Key Properties}

\begin{itemize}
    \item Measures:
    \begin{itemize}
        \item Strength
        \item Direction of linear relationship
    \end{itemize}
    \item It is:
    \begin{itemize}
        \item \textbf{Normalized covariance}
        \item \textbf{Unit-free}
    \end{itemize}
\end{itemize}

\section{Covariance Example 2 (Continuous Case)}

\textbf{Given:} \\
Joint PDF $f(x,y)$

\textbf{Required:} \\
Compute required quantities (expectations, covariance)

\textbf{Solution Flow (as per slides)}

\begin{enumerate}
    \item Start from joint PDF
    \item Compute marginal densities
    \item Compute expectations:
    \[
    E[X], \quad E[Y]
    \]
    \item Compute:
    \[
    E[XY] = \int \int xy f(x,y) \, dx,dy
    \]
    \item Substitute into:
    \[
    \text{Cov}(X,Y) = E[XY] - E[X]E[Y]
    \]
    \item Evaluate integrals step-by-step
\end{enumerate}

\textbf{Final:} \\
Covariance obtained after full substitution and simplification (as shown across slides).

\section{Joint Gaussian Random Variables}

\subsection*{Definition}

A pair $(X, Y)$ is \textbf{jointly Gaussian} if:
\begin{itemize}
    \item Joint density has \textbf{bivariate Gaussian form}
    \item Marginals are \textbf{Gaussian}
\end{itemize}

\subsection*{Interpretation}

\begin{itemize}
    \item Used to model variables affected by many random factors
\end{itemize}

\subsection*{Real-life Example (from slide)}

\begin{itemize}
    \item $(X)$: CO$_2$ concentration
    \item $(Y)$: Temperature
\end{itemize}

Their joint behavior can be approximated by a \textbf{joint Gaussian distribution}.

\subsection*{Applications}

\begin{itemize}
    \item Signal processing
    \item Communication systems
    \item Noisy measurements
\end{itemize}

\subsection*{Key Idea}

Joint Gaussian RVs:
\begin{itemize}
    \item Represent \textbf{continuous random behavior}
    \item Capture dependence through covariance
\end{itemize}

\section*{Final Revision Flow}

\begin{enumerate}
    \item Start with \textbf{Joint PMF $\rightarrow$ Marginal PMF (summation)}
    \item Measure joint behavior using:
    \begin{itemize}
        \item $E[XY]$ (raw correlation)
    \end{itemize}
    \item Improve measure by centering:
    \begin{itemize}
        \item Covariance
    \end{itemize}
    \item Normalize:
    \begin{itemize}
        \item Pearson coefficient
    \end{itemize}
    \item Extend to continuous case:
    \begin{itemize}
        \item Integrals
    \end{itemize}
    \item Model real systems:
    \begin{itemize}
        \item Joint Gaussian RVs
    \end{itemize}
\end{enumerate}

\end{document}